% !TeX root = ../../Book3.tex
\chapter{Poisson–Jensen 公式与零点分布}
\label{b3:ch:18}
\BookChapterNumber{b3:ch:18}
\begin{chapterabstract}
零点的重数决定对数模的点源质量，零点的位置决定它对圆周平均的贡献。本章先建立整数计数、积分计数及其在全纯扰动下的稳定性，再从圆周对数核和 Green 核推导 Jensen 与 Poisson--Jensen 公式。中心零点、观察圆周上的零点及亚纯函数的极点分别处理。随后以增长界控制零点密度，证明有界圆盘全纯函数的零点判据及 Blaschke 乘积构造，最后把这些公式放入次调和函数的 Riesz 测度框架，并证明非零全纯极限下对数模的局部可积收敛。
\end{chapterabstract}

\begin{learninggoals}
\begin{itemize}
\item 按重数记录零点，区分开、闭圆盘计数与积分计数，并用 Rouch\'e 定理追踪扰动中的零点簇。
\item 推导 Jensen 与 Poisson--Jensen 公式，核对中心系数、Green 项的符号以及边界对数奇性的可积性。
\item 把最大模或圆周幂平均的增长界转为零点限制，识别两半径估计中半径间隔的作用。
\item 证明 Blaschke 条件的必要性与充分性，控制局部乘积收敛及除去零点后的有界余因子。
\item 由对数模的 Riesz 测度解释径向平均和零点稳定性，区分全纯函数的正零点测度与亚纯函数的符号零极点测度。
\end{itemize}
\end{learninggoals}

若 \(f(0)=1\)，且 \(|f|\) 在半径 \(R\) 的圆周上不超过十，那么较小圆盘 \(|z|\le r<R\) 中不能放入任意多的零点。一个位于 \(a\) 的零点对平均对数贡献 \(\log(R/|a|)\)，所以较小圆盘中的每个零点至少贡献 \(\log(R/r)\)。要把这个想法变成可靠估计，既要记录重数，也要说明圆周上的对数如何积分；若原点本身是零点，还须先取出它的幂次，不能直接写下 \(\log|f(0)|\)。

\ProseParagraphBreak
\secref{b3:sec:1801}先把零点整理为局部有限的整数质量。整数计数随半径跳跃，积分计数则给每个非中心零点加上对数权重，在对数半径上连续而分段线性。Rouch\'e 定理保持一个小围圈中的总重数，缩小这些围圈便能追踪合并或分裂的零点簇。这种内部稳定性容许质量移向边界，但要求极限函数不恒为零。

\ProseParagraphBreak
\secref{b3:sec:1802}把计数与函数大小联系起来。先算一个对数核的圆周平均，再有限分解零点，便得到含中心修正的 Jensen 公式。任意内部观察点需要用 Green 核替代中心的对数权重；核的负 Laplace 与 \(\log|f|\) 的正点质量恰好抵消。边界零点仍可处理，但必须保留可积对数奇性，不能把边界无零点版本原样套用。两半径估计和 Blaschke 条件随后说明：增长限制不仅约束零点数量，也约束它们接近边界的速度。

\ProseParagraphBreak
\secref{b3:sec:1803}说明这些计算为何属于势论。零点对应正 Riesz 质量，极点则对应负质量；一般次调和函数的质量还可以分布在曲线或面积上。局部一致收敛的全纯函数在移动零点附近未必有一致收敛的对数，但有限零点分解与对数核的可积性足以给出局部 \(L^1\) 收敛。阅读这部分时应分别留意点态值、面积积分和测试函数所记录的信息。

\ProseParagraphBreak
本章使用前面的局部零点分解、辐角原理与 Rouch\'e 定理，以及\chapref{b3:ch:15}的次调和函数和 Riesz 测度、\chapref{b3:ch:17}的正 Green 核。有关整函数增长与零点分布的进一步阅读，可参考 Levin 的《Lectures on Entire Functions》\cite{Levin1996Lectures}。本章的平均公式与有界零点构造在正文中独立建立，后续整函数乘积理论将继续处理给定零点与全局增长之间的关系。

% !TeX root = ../../Book3.tex
\section{零点计数}
\label{b3:sec:1801}

辐角原理用一条闭曲线计算零点总数，但没有直接显示零点到圆心的距离。对数模的圆周平均会把这项几何信息补进计数。

\subsection{重数}
\label{b3:subsec:180101}

非恒零全纯函数在零点 \(a\) 可唯一写成 \(f(z)=(z-a)^m g(z)\)，其中 \(g(a)\ne0\)，\(m\) 是零点重数。后面的计数均按重数重复零点。紧子集上只有有限个零点；否则有内部聚点，恒等定理迫使函数恒零。

例如 \(f(z)=z^2(z-1)^3\) 在原点和一的质量分别为二、三。零点位置相同而重数不同，会改变对数模的奇性强度；按不同零点个数而不计重数不能恢复解析信息。

\subsection{计数函数}
\label{b3:subsec:180102}

对 \(f(0)\ne0\)，令 \(n(r)\) 为 \(|a|\le r\) 内按重数的零点数，定义积分计数
\[
N(r)=\int_0^r\frac{n(t)}t\,\mathrm{d}t
=\sum_{|a|<r}\log\frac r{|a|}.
\]
第二式由有限和换序得到。零点在半径边界时贡献零，故开闭圆盘约定不影响 \(N\)。若原点重数为 \(m\)，先除去 \(z^m\)，再单独保留 \(m\log r\) 项。\(n\) 是阶梯函数，\(N\) 却连续；每个零点从它进入圆盘时开始累计对数权重。

\subsection{辐角原理再述}
\label{b3:subsec:180103}

边界无零点时，\cref{b3:thm:03-argument-principle}给
\[
n(r)=\frac1{2\pi\mathrm{i}}\int_{|z|=r}\frac{f'(z)}{f(z)}\,\mathrm{d}z.
\]
这里要求 \(f\) 在闭圆盘邻域全纯，或者先在稍小半径应用；边界上有零点时积分不定义。辐角原理显示整数计数的拓扑稳定性，积分计数则更适合平均与增长估计。后面不会把 \(n(r)\) 的跳跃误当作对数平均本身的跳跃。


\ProseParagraphBreak

计数函数的两个版本表达不同信息。若半径内有 \(N\) 个零点全位于 \(|a|=r/2\)，它们给 \(n(r)=N\)、\(N(r)=N\log2\)；若相同个数全位于 \(|a|=(1-\varepsilon)r\)，整数计数仍为 \(N\)，积分计数却为 \(-N\log(1-\varepsilon)\)，随 \(\varepsilon\downarrow0\) 趋零。紧贴观察圆周的零点对该半径的平均影响较小，但一旦把半径扩大，它们会继续累积权重。因此从 \(N(r)\) 控制 \(n(r)\) 必须另取更大的半径，不能同半径直接比较而没有额外系数。

\ProseParagraphBreak

换一个圆心可以避开原点零点。对内部点 \(b\) 且 \(f(b)\ne0\)，在 \(B(b,R)\) 使用变量 \(\zeta=z-b\)，得到完全相同的计数和公式。如果要在几个不同圆心比较，应分别保留 \(\log|f(b)|\)；这个中心常数可能很负，会使零点上界较弱。例如函数在 \(b\) 非零但已非常接近一个零点时，\(|f(b)|\) 小，计数估计仍成立却不一定有效。归一化 \(f(b)=1\) 可消除常数项，但乘法归一化并不改变零点本身。

\ProseParagraphBreak

对含原点重数的函数，用修正积分计数
\[
N_f(r)=m\log r+\sum_{0<|a|<r}m_a\log(r/|a|)
\]
十分方便。它在小于一的半径可能为负，因为 \(m\log r<0\)；这不是非原点零点贡献出现负值，而是已单独记录的中心奇性。若只需要非负零点数的估计，先除去原点幂再应用公式更稳妥。所有符号版本必须在中心奇性项上保持同一约定，尤其进入亚纯函数的零点、极点差计数时不能重复加上原点项。

\ProseParagraphBreak

零点对全纯扰动的稳定性有一个比总数更精细的表述。设 \(f_j\to f\) 在区域 \(D\) 内局部一致，且 \(f\not\equiv0\)。固定一个内部零点 \(a\)，取小圆盘使闭盘中没有其他极限零点，边界上的 \(|f|\) 有正下界。Rouch\'e 定理说明：充分大的 \(j\) 在盘内恰有 \(m_a\) 个零点，按重数计。把圆盘继续缩小，仍有同样结论，所以这些零点全向 \(a\) 集中。紧集上远离极限零点的部分，由 \(|f|\) 的正下界可知最终没有新的零点。
\ProseParagraphBreak
于是离散测度 \(\mu_j=\sum m_{j,a}\delta_a\) 在内部弱收敛到 \(\mu=\sum m_a\delta_a\)：对紧支集连续测试函数，只需用有限个互不相交小圆盘包住相关极限零点，每盘总质量相同，测试函数在盘内的振幅任意小，其余部分最终无质量。这里弱收敛只针对内部紧支集；零点可以逃向边界或无穷远。例如 \(f_j(z)=1-z/j\) 在全平面局部一致趋于一，唯一零点趋向无穷远。极限函数没有零点，与每个近似函数有一个零点并不矛盾。
\ProseParagraphBreak
假设极限不恒零是必要的。\(f_j(z)=e^{-j}\sin(jz)\) 在足够窄的固定水平带中局部一致趋于零，但零点间距趋于零，任何固定实区间内的零点数都趋向无穷。恒等为零的函数不再拥有离散的零点测度，不能给它随意指定一个有限重数来延续上面的结论。

% !TeX root = ../../Book3.tex
\section{Jensen 公式}
\label{b3:sec:1802}

除掉圆盘内的零点，剩余函数的对数模调和。中心平均等式因此会留下每个零点的一个对数距离项，这就是 Jensen 公式。

\subsection{Jensen 公式}
\label{b3:subsec:180201}

\term[jensen-gongshi]{Jensen 公式}[Jensen's formula]的归一化以圆周平均 \(\mathrm{d}t/(2\pi)\) 为准。
\begin{theorem}[Jensen]\label{b3:thm:18-jensen}
设 \(f\) 在闭圆盘 \(\overline{B(0,R)}\) 邻域全纯，\(f(0)\ne0\)，边界无零点。内零点按重数记为 \(a_1,\ldots,a_N\)。则
\[
\frac1{2\pi}\int_0^{2\pi}\log|f(Re^{\mathrm{i}t})|\,\mathrm{d}t
=\log|f(0)|+\sum_{j=1}^N\log\frac R{|a_j|}.
\]
\end{theorem}
\begin{proof}
令 \(g=f/\prod(z-a_j)\)，它在闭圆盘无零点，故 \(\log|g|\) 调和。每个因子满足圆周平均 \(\log|Re^{\mathrm{i}t}-a|=\log R\)：写成 \(\log R+\operatorname{Re}\log(1-(a/R)e^{-\mathrm{i}t})\)，收敛级数的每个非零角频率平均零。于是边界平均为 \(\log|g(0)|+N\log R\)，而 \(\log|g(0)|=\log|f(0)|-\sum\log|a_j|\)，得公式。
\end{proof}
边界零点的对数奇性可积；用相邻半径与局部零点分解可把公式延伸到这样的半径。原点为零点时，用 \(f=z^m g\)、\(g(0)\ne0\)，中心常数改成 \(\log|g(0)|\)，并增加 \(m\log R\)。

\subsection{Poisson–Jensen 公式}
\label{b3:subsec:180202}

\begin{theorem}[Poisson--Jensen]\label{b3:thm:18-poisson-jensen}
在上述条件下，对非零点 \(z\in B(0,R)\)，有
\[
\log|f(z)|=\frac1{2\pi}\int_0^{2\pi}P_R(z,t)\log|f(Re^{\mathrm{i}t})|\,\mathrm{d}t
-\sum_jg_{B(0,R)}(z,a_j).
\]
\end{theorem}
\begin{proof}
加上各个正 Green 核后，函数 \(\log|f|+\sum_jg(\cdot,a_j)\) 在每个零点的 \(m\log|z-a_j|\) 奇性完全抵消，成为调和函数。Green 核边界零，所以边值仍为 \(\log|f|\)。圆盘 Poisson 唯一性给表示。零点处两侧均按极限解释为负无穷。
\end{proof}
亚纯函数的版本在每个极点添加对应的正 Green 项：\(\log|f|\) 的 Laplace 测度为零点质量减极点质量。中心公式因此是零点积分计数减极点积分计数。

\subsection{零点分布估计}
\label{b3:subsec:180203}

\begin{corollary}\label{b3:cor:18-zero-estimate}
若 \(f\) 在 \(B(0,R)\) 全纯，\(f(0)\ne0\)，\(M(S)=\max_{|z|=S}|f(z)|\)，则 \(0<r<S<R\) 时
\[
n(r)\log\frac Sr\le\log M(S)-\log|f(0)|.
\]
\end{corollary}
\begin{proof}
半径 \(r\) 内每个零点向半径 \(S\) 的 Jensen 和贡献至少 \(\log(S/r)\)。平均对数不超过最大对数，代入公式即得。边界零点用稍移半径处理。
\end{proof}
圆盘上有界全纯函数的零点满足 Blaschke 条件 \(\sum(1-|a_j|)<\infty\)：令 \(S\uparrow1\)，Jensen 和有界，再用 \(-\log t\ge1-t\)。有界性因而不允许零点过快地向整个圆周拥挤。


\ProseParagraphBreak

边界零点的延伸可写得更明确。设在 \(|a|=R\) 有一个因子 \((z-a)^m\)，余函数在该点不为零。边界局部参数 \(t=t_0+h\) 时，\(|Re^{\mathrm{i}t}-a|\) 与 \(|h|\) 相比，故 \(\log\) 奇性由常数加 \(m\log|h|\) 控制，而 \(\int_{-\delta}^{\delta}|\log|h||\,\mathrm{d}h<\infty\)。其平均仍为 \(\log R\)：先在 \(|a|<R\) 得到等式，再使 \(a\) 沿半径趋向边界，局部负部的积分可由该对数估计一致控制。Jensen 和中该零点贡献 \(\log(R/|a|)=0\)，因此公式在边界零点半径保持正确。这里需要的是对数的积分收敛，不能只凭逐点极限使用有界收敛定理。

\ProseParagraphBreak

一阶函数已把公式的几何意义展示出来。对 \(f(z)=1-az\)、\(a\ne0\)，有
\[
\frac1{2\pi}\int_0^{2\pi}\log|1-aRe^{\mathrm{i}t}|\,\mathrm{d}t
=\max(0,\log(|a|R)).
\]
当 \(|a|R<1\) 时，圆盘无零点，平均等于中心对数零；当 \(|a|R>1\) 时，一个零点进入圆盘，平均变成 \(\log(|a|R)\)。临界半径仍连续，只有对 \(\log R\) 的斜率发生跳跃。更一般地，对多项式 \(c\prod(z-a_j)\)，边界平均为 \(\log|c|+\sum_j\log\max(R,|a_j|)\)。这可用于不求全部积分而从根的位置直接计算平均。

\ProseParagraphBreak

Jensen 公式控制的是平均对数，不是模的算术平均。由对数凹性，
\[
 \exp\left(\frac1{2\pi}\int_0^{2\pi}\log|f(re^{\mathrm{i}t})|\,\mathrm{d}t\right)
 \le\frac1{2\pi}\int_0^{2\pi}|f(re^{\mathrm{i}t})|\,\mathrm{d}t.
\]
左侧是几何平均。边界上可能有很小的函数值，也可能有很大的值；几何平均仍把所有内部零点贡献准确相加。例如 \(f(z)=e^{az}\) 没有零点，其对数平均为零，而模的平均一般大于一。零点资料完全看不见指数整函数中的调和对数模，正好与上一章 Riesz 测度的自由度相符。

\ProseParagraphBreak

Poisson--Jensen 公式的单个零点项也能显式核验。在单位圆盘，取 \(f(z)=z-a\)，\(|a|<1\)，则
\[
\log|z-a|+g_{\mathbb D}(z,a)=\log|1-\overline az|.
\]
右边无零点且调和，边界上两项对数模相等。因而 Poisson 延拓的是 \(\log|1-\overline az|\)，减去 Green 核后才恢复真正的 \(\log|z-a|\)。若直接对边界 \(\log|f|\) 作 Poisson 积分并宣称得到内部 \(\log|f|\)，就遗漏了零点的负对数奇性。错误在内部零点附近最明显：Poisson 解有限，原对数模却趋负无穷。

\ProseParagraphBreak

亚纯版本可通过有限因子分解完整证明。设 \(f\) 在闭圆盘邻域亚纯，边界无零点、极点，中心既不是零点也不是极点。内零点为 \(a_j\)、极点为 \(b_k\)，均按重数重复。将 \(f\) 乘 \(\prod(z-b_k)\) 再除 \(\prod(z-a_j)\)，所得函数在闭圆盘无零点全纯。对它应用调和平均，分别计算因子平均，得到
\[
\frac1{2\pi}\int\log|f(Re^{\mathrm{i}t})|\,\mathrm{d}t
=\log|f(0)|+\sum_j\log(R/|a_j|)-\sum_k\log(R/|b_k|).
\]
对一般内部点同样加零点 Green 核、减极点 Green 核，奇性抵消后使用 Poisson 表示即可。中心为零点或极点时，先提出整数次幂，指数允许为负。这个公式精确记录零点与极点的差，但它们数量同时很大时可能抵消，所以并不单独控制两类总量。

\ProseParagraphBreak

零点估计中选取半径比会决定常数。取 \(S=2r\) 得 \(n(r)\le(\log M(2r)-\log|f(0)|)/\log2\)；若取 \(S=(1+\varepsilon)r\)，分母 \(\log(1+\varepsilon)\) 很小，估计通常变弱。若 \(f\) 只在固定圆盘全纯，不能为了得到方便的二倍半径超出定义域。靠近外边界时，应选仍在域中的 \(S\)，并保留其实际系数。

\ProseParagraphBreak

作为精确的禁止条件，假设单位圆盘有界非恒零全纯函数的零点为 \(a_n=1-1/n\)、\(n\ge2\)。因为 \(\sum(1-|a_n|)=\sum1/n\) 发散，这种函数不存在；但没有有界性限制时，该点列在圆盘内部无聚点，仍可通过区域上的零点构造实现。当前整平面 Weierstrass 定理尚不能直接用于这类边界聚点，须使用区域版本，或者取 \(\sin(\pi/(1-z))\) 在圆盘上的明确函数，它确有这些零点并在接近边界时无界。这把解析存在性与有界函数的增长限制分开。


\ProseParagraphBreak

有限零点还能给逐点估计。对 \(a\in\mathbb D\setminus\{0\}\) 定义
\[
b_a(z)=\frac{|a|}{a}\frac{a-z}{1-\overline az},
\qquad b_0(z)=z.
\]
它在 \(a\) 有单零点，圆周模一，圆盘内部模小于一。若 \(|f|\le1\)，将任意有限批零点提出，\(F=f/\prod b_{a_j}\) 在圆盘全纯。在 \(|z|=r\uparrow1\) 上，有限乘积的模一致趋一，故最大模原理给 \(|F(z)|\le1\)。因此
\[
|f(z)|\le\prod_j|b_{a_j}(z)|.
\]
原点处即 \(|f(0)|\le\prod|a_j|\)，与 Jensen 完全一致；在零点以外的其他观察点，这个估计保留了角位置，通常比只按半径计数的信息更多。若原点是零点，则先提出原点因子，不能在原点两边同时为零的平凡不等式中读取其他零点数量。

\ProseParagraphBreak

对于任意离散零点列，Blaschke 条件在单位圆盘也是有界实现的充分条件。这可以在当前章独立构造，毋须将下一章的整平面乘积定理倒用。设 \(\sum(1-|a_j|)<\infty\)，只需先分离原点的有限重数。写 \(a=re^{\mathrm{i}\theta}\)，则
\[
1-b_a(z)=(1-r)\frac{1+e^{-\mathrm{i}\theta}z}{1-r e^{-\mathrm{i}\theta}z},
\quad
\sup_{|z|\le s}|1-b_a(z)|\le\frac{1+s}{1-s}(1-r),\ s<1.
\]
尾部因此在每个内部紧集上有可和误差。取误差充分小的尾因子，用原点附近的对数级数 \(\log(1+w)\)，其模不超过 \(2|w|\)，所以对数和局部一致收敛，指数给非零尾乘积。有限头部保留准确零点及重数，所有部分乘积模不超过一，极限也是有界整于圆盘的函数。乘积在原点为 \(\prod|a_j|>0\)：从足够大的 \(j\) 起，\(-\log|a_j|\le2(1-|a_j|)\)，故对数和有限。这证明极限不恒零。必要性与充分性由此在当前范围完整配对。

\ProseParagraphBreak

注意该乘积是单位圆盘全纯函数，一般不能延拓为整个复平面的整函数。其分母的零点在 \(1/\overline a_j\)，可能向圆周聚集；构造只在内部紧集上收敛。称为 Blaschke 乘积时，应把定义域和模界一起记住。条件 \(\sum(1-|a_j|)<\infty\) 也不保证普通边界每一点的连续性，径向边界理论将由后面的函数空间条件进一步研究。

\ProseParagraphBreak

例如 \(a_j=1-2^{-j}\) 满足条件，给一个无穷零点且有界的圆盘全纯函数。点列只有边界聚点一，故不违反恒等定理。若改成 \(a_j=1-j^{-1}\)，条件失去，可在圆盘全纯但必须无界的函数中实现。若将每个 \(1-2^{-j}\) 重复 \(2^j\) 次，按不同零点看仍是原点列，但带重数的缺陷和变成 \(\sum1\)，再次不能有界实现。重数在有界零点判据中的作用不亚于位置。

\ProseParagraphBreak

即使没有统一最大模上界，圆周上的模积分也能控制零点。若某个 \(p>0\) 满足
\[
\sup_{0<r<1}\frac1{2\pi}\int_0^{2\pi}|f(re^{\mathrm{i}t})|^p\,\mathrm{d}t\le A<\infty,
\]
对数凹性给 \((2\pi)^{-1}\int\log|f|\le p^{-1}\log A\)，先以正截断避免零点，再取极限。Jensen 的零点和仍有统一上界，所以非恒零函数继续满足 Blaschke 条件。这一推导解释为何后面的 \(H^p\) 空间即使含有无界函数，零点仍受同一个可和性限制。此处已给所需零点论证，不把后面的空间名称当成未建立的前置。

\ProseParagraphBreak

假设中 \(p>0\) 不能简单换成“平均 \(\log|f|\) 有有限值”：每个固定内部圆周的平均对数本来就是有限的，只有对所有半径的共同上界才能控制边界零点和。量词的差别在 \(\sin(\pi/(1-z))\) 上可看出，它在每个内部圆周的平均都有限，但随着半径逼近一，越来越多零点进入，累计对数和无界。逐个半径的良定义不等于统一增长控制。

\ProseParagraphBreak

有界函数的零点分解还可以连同余因子的模界一起得到。设 \(f\not\equiv0\)、\(|f|\le1\)，取含全部零点的 Blaschke 乘积 \(B\)，包括原点的有限幂。每个有限部分乘积除出后，前面已经证明所得函数的模不超过一。在不含零点的紧集上令部分乘积趋于 \(B\)，便有 \(|f/B|\le1\)；在共同零点处，重数恰好相同，所以商的奇点可去。故
\[
 f=B F,\qquad F\text{ 在圆盘全纯、无零点，且 }|F|\le1.
\]
无零点并不意味着常数。正调和函数 \(-\log|F|\) 记录全部零点之后仍剩下的自由度。
\ProseParagraphBreak
若再设 \(f(0)\ne0\)，Harnack 不等式可把中心值变成内部下界。由于 \(|B(0)|\le1\)，有 \(|F(0)|\ge|f(0)|\)。对 \(|z|\le r<1\)，令 \(v=-\log|F|\ge0\)，便得
\[
 |F(z)|=e^{-v(z)}\ge |F(0)|^{(1+r)/(1-r)}
                  \ge |f(0)|^{(1+r)/(1-r)}.
\]
因此 \(|f(z)|\ge|B(z)|\,|f(0)|^{(1+r)/(1-r)}\)。这个估计保留了每个零点对观察点的伪双曲距离；若 \(z\) 靠近零点，\(|B(z)|\) 自然趋零，估计不会错误地宣称函数处处远离零。若已知某个紧集与零点保持分离，再估计其中的乘积，才能得到实际的正下界。
\ProseParagraphBreak
一个特别有启发性的无零点余因子是
\[
 S(z)=\exp\left(-\frac{1+z}{1-z}\right),\qquad
 -\log|S(z)|=\frac{1-|z|^2}{|1-z|^2}.
\]
右端是边界点1的 Poisson 核，所以在圆盘正，\(|S|<1\)，且无零点。对每个 \(e^{\mathrm{i}\theta}\ne1\)，径向边界模趋于一；沿正实轴趋向1时，函数模却趋于零。Jensen 平均始终为 \(\log|S(0)|=-1\)。如果只把几乎处处的边界模一代入平均，就会得到零，丢失了集中到单个边界点的质量。这说明内部 Poisson 表示与边界极限交换需要可积控制，也说明零点资料不能决定一个有界全纯函数。

% !TeX root = ../../Book3.tex
\section{次调和函数观点}
\label{b3:sec:1803}

Jensen 公式还可以读成一个势论恒等式：圆周平均之差等于圆盘内部 Laplace 测度的对数势。零点只是这一原则的离散实例。

\subsection{\texorpdfstring{\(\log |f|\)}{log|f|}}
\label{b3:subsec:180301}

函数 \(\log|f|\) 在零点外调和，零点处局部为 \(m\log|z-a|\) 加调和函数。这个分解不依赖对数分支选择，因为只取实部。它也使每个局部可积性与奇性符号都可以直接核验。

\subsection{作为 Riesz 测度的零点}
\label{b3:subsec:180302}

\begin{proposition}\label{b3:prop:18-zero-measure}
区域上非恒零全纯函数满足
\[
\frac1{2\pi}\Delta\log|f|=\sum_am_a\,\delta_a,
\]
其中右侧记号 \(m_a\) 表示 \(f\) 在 \(a\) 的零点重数。
\end{proposition}
\begin{proof}
在非零点邻域 Laplace 为零；在零点邻域用局部分解与对数基本解，得到重数乘点质量。零点局部有限，因此各局部恒等式通过有限划分拼成分布等式。
\end{proof}
对亚纯函数，把极点计为负重数，即得相应符号测度。

\subsection{势论解释}
\label{b3:subsec:180303}

若 \(u\in C^2\) 次调和，中心圆周平均 \(m(r)\) 的导数为 \(\mu_u(B(0,r))/r\)，所以
\[
m(R)-u(0)=\int_0^R\frac{\mu_u(B(0,t))}t\,\mathrm{d}t.
\]
证明只需积分圆周平均导数，归一化由 \(\mu_u=\Delta u/(2\pi)\) 确定。光滑逼近可把这一公式推广到适当的次调和代表与非奇中心。对 \(u=\log|f|\)，测度离散，右端正是零点对数计数，因而回到 Jensen。

中心为对数奇点时，\(u(0)=-\infty\)，不能直接相减；必须先分离中心点质量。这是原点重数项出现的原因。

\subsection{与 Nevanlinna 理论的联系}
\label{b3:subsec:180304}

对于亚纯函数，零点与极点可能同时增长。单独看 \(\log|f|\) 的平均会发生抵消，而 \(\log^+|f|\) 的平均与极点积分计数之和能够保留总增长信息。这引向\appref{b3:app:C}的 Nevanlinna 特征。

目前已经得到的是零点的精确势论记录，并没有从 Jensen 一式直接获得第二基本定理。后者还需要对数导数和多个值之间的比较。


\ProseParagraphBreak

一般次调和平均公式可借上一章的局部 Riesz 分解在书内推出。取较大内部圆盘，使 \(u=V+h\)，\(V(z)=\int\log|z-w|\,\mathrm{d}\mu(w)\)，\(h\) 调和。对固定圆周的核平均，有
\[
\frac1{2\pi}\int_0^{2\pi}\log|re^{\mathrm{i}t}-w|\,\mathrm{d}t=\log\max(r,|w|).
\]
若 \(u(0)>-\infty\)，势在中心亦有限，故减去中心值后，核贡献 \(\log^+(r/|w|)\)，这是非负函数。Tonelli 给
\[
\frac1{2\pi}\int u(re^{\mathrm{i}t})\,\mathrm{d}t-u(0)
=\int_{|w|<r}\log(r/|w|)\,\mathrm{d}\mu(w)
=\int_0^r\frac{\mu(B(0,t))}{t}\,\mathrm{d}t.
\]
最后一步将 \(\log(r/|w|)\) 写成 \(\int_{|w|}^r\mathrm{d}t/t\)，再换非负积分次序。圆周上原势的负部可积由此也得到控制；无需未经说明地从分布等式取点值。

\ProseParagraphBreak

这个平均差随 \(r\) 非减，因为 Riesz 测度为正。在对数半径 \(s=\log r\) 下，平均函数的导数为 \(\mu(B(0,e^s))\)，它也非减，所以平均关于 \(s\) 凸。对全纯对数模，导数就是零点整数计数；对一般次调和函数，导数可以连续变化。取 \(u(z)=|z|^2\) 时圆周平均为 \(r^2\)，\(\mu(B(0,r))=2r^2\)，导数为 \(2r\)，恰与公式 \(\mu(B(0,r))/r\) 一致。这同时核验面积归一化和 \(2\pi\) 常数。

\ProseParagraphBreak

若中心是零点，\(u(0)=-\infty\)，则两边不能写作有限数相减。对 \(u=m\log|z|+v\)、\(v(0)\) 有限，圆周平均先保留 \(m\log r\)，再对 \(v\) 使用上述公式；测度中先去掉 \(m\delta_0\)。若任意在原点补一个有限值，Riesz 分布不会因单点修改改变，但次调和平均不等式会失败。这说明分布等价类与点态次调和代表应保持区别，零点处的负无穷值具有数学作用。

\ProseParagraphBreak

对于亚纯函数的 \(\log|f|\)，极点给正无穷的对数奇性，函数本身一般既非次调和也非超调和；它的 Laplace 是符号测度。可以局部写成两个次调和函数之差，对零、极点分别应用平均公式，再取差。符号测度下积分换序需要有限性或正负分解，不能继续把整个表达当成非负积分使用 Tonelli。将正测度论证中的每一步直接复制给符号测度，是从全纯版本过渡到亚纯版本时常见的漏洞。

\ProseParagraphBreak

Nevanlinna 特征保留大小的正部及极点总计数。此处先作一个算术检验：令 \(f(z)=1/z\)，对 \(r>1\)，圆周上 \(\log^+|f|=0\)，但中心极点计数贡献 \(\log r\)，所以特征仍增长；对 \(f(z)=z\)，没有极点，正部平均为 \(\log r\)。两函数互为倒数，粗略增长规模相同，单独的平均 \(\log|f|\) 却一正一负。正部与极点项的结合正是为了防止这种符号抵消丢失增长资料。完整定义及中心归一化将在相应附录中给出，这里只解释为何需要新的量。

\ProseParagraphBreak

零点测度的稳定性也可以从对数模的局部可积收敛看出。下面补上这项分析事实。仍设 \(f_j\to f\not\equiv0\) 局部一致，选一个闭圆盘，其边界无极限零点，内部零点有限。由前面的稳定计数，令 \(p_j\) 为盘中全部零点构成的首一多项式，\(p\) 为极限零点多项式，则 \(p_j\to p\) 一致。商 \(h_j=f_j/p_j\) 在盘内全纯无零点，边界上一致趋于 \(h=f/p\)；最大模原理把收敛传到闭盘。\(h\) 无零点且闭盘上有正下界，故 \(\log|h_j|\to\log|h|\) 一致。
\ProseParagraphBreak
剩下只需处理移动的对数核。若 \(a_j\to a\)，则 \(\log|z-a_j|\to\log|z-a|\) 在固定紧集上依 \(L^1\) 收敛。证明时先删去 \(a\) 附近半径 \(2\delta\) 的小圆，在其外一致收敛；其内，所有充分大的 \(a_j\) 落在半径 \(\delta\) 内，负部积分由
\(2\pi\int_0^{3\delta}r|\log r|\,\mathrm{d}r\) 控制，趋于零。正部在固定有界区域统一有界，所删小圆的面积也趋零。两部分合并即得结论。把有限个移动零点项相加，便得到
\[
 \log|f_j|\longrightarrow\log|f|\quad\text{于 }L^1_{\mathrm{loc}}(D).
\]
对紧支集光滑测试函数 \(\varphi\)，积分 \(\int\log|f_j|\Delta\varphi\) 遂收敛；由零点的 Riesz 测度公式，正好再次得到零点测度的弱收敛。点态对数在移动零点处可能完全没有共同下界，局部可积性才是传递 Laplace 算子的合适层次。
\ProseParagraphBreak
这项结论也说明为什么在正规族中控制函数本身，经常足以控制零点，但前提是先排除恒零极限。若极限恒零，\(\log|f_j|\) 可能整体趋于负无穷，即使零点数并未增加；例如常数 \(1/j\) 的对数在每个点都趋于负无穷，而其 Laplace 始终为零。对数的加性常数与零点测度是两种独立资料，不能从测度收敛反推对数模已经收敛。


\begin{chaptersummary}
零点计数有三个相互联系的表达：辐角积分给出围道内的总重数，积分计数给每个零点赋予径向对数权重，对数模的 Laplace 则把重数写成点质量。整数计数在观察圆周经过零点时跳跃，积分计数与圆周平均保持连续；在对数半径上，斜率的跳跃正好记录该圆周上的质量。

\ProseParagraphBreak
Jensen 公式由圆周对数核平均和有限零点分解得到，中心有零点时保留其幂次及首个非零系数。Poisson--Jensen 公式用正 Green 核表示任意内部点的零点影响：零点项从边界调和延拓中减去，极点项则加回。边界零点的对数奇性可积，公式能否使用取决于所取版本及其边界处理，不能仅凭路径经过零点就否定一切平均表达。

\ProseParagraphBreak
函数增长给零点分布加上定量限制。两半径估计利用每个内部零点的共同正权重；在单位圆盘中让外半径趋于一，就得到计入重数的 Blaschke 条件。归一化因子使尾项趋于一，局部对数级数保证乘积收敛并保留精确重数。除去全部零点后仍有一个有界、无零点的全纯余因子，它可以非常数，零点资料因此不足以恢复原函数。

\ProseParagraphBreak
次调和平均公式把离散零点推广为一般正质量，其单侧导数区分开、闭圆盘。非零全纯极限下，Rouch\'e 定理先稳定局部零点簇，再由移动对数核的可积收敛得到对数模的局部 \(L^1\) 稳定性。极限恒零或同时出现正、负零极点质量时，需要重新检查归一化和测试函数范围；原有的正测度结论不能直接照搬。
\end{chaptersummary}

\BookExercises

% 自拟：把圆盘径向平均推广到圆环，区分整数绕数与盘内零点数，并解释允许负斜率的原因。
\Needspace{6\baselineskip}
\begin{exercise}[圆环上的对数平均]\label{b3:ex:18-draft-01}
设 \(0<\rho<R<\infty\)，\(f\) 在圆环 \(A=\{\rho<|z|<R\}\) 全纯且不恒为零。令
\[
 M(t)=\frac1{2\pi}\int_0^{2\pi}\log|f(e^{t+\mathrm i\theta})|\,\mathrm d\theta,
 \quad \log\rho<t<\log R.
\]
证明 \(M\) 连续，在每个紧区间上分段仿射，并在一维分布意义下满足
\[
 M''=\sum_{a\in Z(f)}m_a\delta_{\log|a|}.
\]
在圆周 \(|z|=e^t\) 无零点时，证明
\[
 M'(t)=\frac1{2\pi\mathrm i}\int_{|z|=e^t}\frac{f'(z)}{f(z)}\,\mathrm dz\in\mathbb Z.
\]
由此写出从任一无零点初始圆周 \(t_0\) 到 \(t>t_0\) 的平均增量公式，并举例说明 \(M\) 虽然凸，却未必非减。
\end{exercise}
\begin{solution}
取任意紧区间 \([t_1,t_2]\subset(\log\rho,\log R)\)，在其对应的闭圆环外再留出一点余量。较大的闭圆环中只有有限个零点，除去相应有限多项式因子后，余函数在这段圆环附近全纯无零点。因此
\[
 \log|f(z)|=h(z)+\sum_a m_a\log|z-a|,
\]
其中 \(h\) 调和，和式有限。即使余函数在整个圆环上未必有单值全纯对数，其对数模仍由局部全纯对数的实部给出，是单值调和函数。

令 \(H(t)\) 为 \(h(e^{t+\mathrm i\theta})\) 的角平均。极坐标下的调和方程给
\[
 \partial_t^2h(e^{t+\mathrm i\theta})
 +\partial_\theta^2h(e^{t+\mathrm i\theta})=0.
\]
对 \(\theta\) 积分时，周期性使第二项积分为零，故 \(H''=0\)。由圆周对数平均公式，在 \([t_1,t_2]\) 上有
\[
 M(t)=\alpha t+\beta+\sum_a m_a\max\{t,\log|a|\}.
\]
这同时证明了平均的有限性、连续性和局部分段仿射性质。函数 \(\max\{t,s\}\) 的二阶分布导数为 \(\delta_s\)，于是得到题中的分布等式；它的右端在每个紧区间上只有有限项。

若观察圆周无零点，可以在邻近圆周一致地对对数模求导，得到
\[
 M'(t)=\operatorname{Re}\frac1{2\pi}
 \int_0^{2\pi}\frac{e^{t+\mathrm i\theta}f'(e^{t+\mathrm i\theta})}
 {f(e^{t+\mathrm i\theta})}\,\mathrm d\theta.
\]
右侧正是题中围道积分的实部。该围道积分是闭曲线 \(f(e^{t+\mathrm i\theta})\) 绕原点的整数绕数，因而本身为实数，得到所需等式。这里 \(f\) 未必定义在圆环内孔中，不能把这个整数直接解释为整张圆盘内的零点数。

记初始圆周的整数为 \(k=M'(t_0)\)。每经过一个零点半径，斜率增加该圆周上的总重数，故
\[
 M(t)=M(t_0)+k(t-t_0)
 +\sum_{e^{t_0}<|a|<e^t}m_a\bigl(t-\log|a|\bigr).
\]
把终点圆周上的零点加进和式不改变值，因为它们的权重为零。二阶导数为正测度，说明 \(M\) 凸；但初始整数 \(k\) 可以为负。例如 \(f(z)=1/z\) 在圆环全纯无零点，\(M(t)=-t\)，在整个区间严格递减。
\end{solution}

% 自拟：求固定中心值下开、闭圆盘的最优整数计数，并完成极值构造与等号分类。
\Needspace{6\baselineskip}
\begin{exercise}[最优零点计数与边界]\label{b3:ex:18-draft-02}
固定 \(0<c<1\)、\(0<r<1\)。在所有满足 \(|f|\le1\)、\(f(0)=c\) 的单位圆盘全纯函数中，分别求闭圆盘 \(|z|\le r\) 与开圆盘 \(|z|<r\) 内零点总重数的最大值。令
\[
 T=\frac{\log c}{\log r},
\]
证明两个答案分别为 \(\lfloor T\rfloor\) 与 \(\lceil T\rceil-1\)，并构造达到它们的函数。若闭圆盘中的零点总重数恰等于实数上界 \(T\)，确定全部这样的函数。
\end{exercise}
\begin{solution}
将闭圆盘中零点按重数列为 \(a_1,\ldots,a_N\)，因 \(f(0)>0\)，这些点都非零。取归一化因子
\[
 b_a(z)=\frac{|a|}{a}\frac{a-z}{1-\overline a z},\quad
 B(z)=\prod_{j=1}^Nb_{a_j}(z).
\]
有限 Blaschke 除法给出 \(f=Bh\)、\(|h|\le1\)，从而
\[
 c=|f(0)|\le B(0)=\prod_{j=1}^N|a_j|\le r^N.
\]
由于 \(\log r<0\)，这意味着 \(N\le T\)，故闭圆盘计数至多 \(\lfloor T\rfloor\)。若所列零点都在开圆盘且 \(N\ge1\)，第二个不等式严格，故 \(N<T\)；\(N=0\) 时同样有 \(N<T\)，因为 \(T>0\)。开圆盘中的整数上界因而为 \(\lceil T\rceil-1\)。

为达到闭圆盘上界，令 \(N=\lfloor T\rfloor\)。若 \(N=0\)，取 \(f=c\)。若 \(N\ge1\)，任选 \(N\) 个模为 \(r\) 的点，允许按重数重复，以它们组成 \(B\)，并取
\[
 f(z)=\frac{c}{r^N}B(z).
\]
这里 \(c\le r^N\)，所以 \(|f|\le1\)，且 \(f(0)=c\)，零点恰为选定的 \(N\) 个。

对于开圆盘，令 \(N=\lceil T\rceil-1\)。\(N=0\) 时仍取常函数。若 \(N\ge1\)，则 \(N<T\) 给出 \(c<r^N\)，可以选择 \(c^{1/N}\le s<r\)。将上面的构造改成 \(N\) 个模为 \(s\) 的零点及系数 \(c/s^N\)，即可达到这个上界。

最后，若闭圆盘计数为 \(N=T\)，则 \(T\) 必须是正整数，且 \(c=r^N\)。前面的整条不等式都取等号，故每个 \(|a_j|=r\)，并有 \(|h(0)|=1\)。最大模原理使 \(h\) 为模一常数；又因归一化后 \(B(0)=r^N=c=f(0)\)，这个常数只能为一。因此全部极值函数正是由 \(N\) 个模为 \(r\) 的零点组成的归一化有限 Blaschke 乘积，不含其他零点。
\end{solution}

% 对原自拟题实质改编：以双重对数修正检验 Blaschke 可和性的两个临界指数。
\Needspace{6\baselineskip}
\begin{exercise}[重数的临界可和性]\label{b3:ex:18-draft-03}
给定 \(\beta,\gamma\in\mathbb R\)，对 \(n\ge3\) 令
\[
 a_n=1-\frac1{n^2},\quad
 m_n=\max\left\{1,
 \left\lfloor\frac{n}{(\log n)^\beta(\log\log n)^\gamma}\right\rfloor\right\}.
\]
确定哪些 \((\beta,\gamma)\) 使这些点以指定重数恰为某个非恒零有界圆盘全纯函数的全部零点。解答须同时给出必要性与存在性。
\end{exercise}
\begin{solution}
对任意固定实数 \(\beta,\gamma\)，有
\[
 x_n=\frac{n}{(\log n)^\beta(\log\log n)^\gamma}\longrightarrow\infty.
\]
这是因为 \(\log x_n=\log n-\beta\log\log n-\gamma\log\log\log n\)，后两项都为 \(o(\log n)\)。因此取整与最大值只改变有限开头及渐近相对误差，\(m_n/x_n\to1\)。带重数的 Blaschke 和与下列级数有相同敛散性：
\[
 \sum_{n\ge3}m_n(1-a_n)
 \quad\text{与}\quad
 \sum_{n\ge3}\frac1{n(\log n)^\beta(\log\log n)^\gamma}.
\]
后一被积型函数最终递减，积分判别法和代换 \(u=\log x\) 将问题化为
\[
 \int_{\log3}^\infty\frac{\mathrm du}{u^\beta(\log u)^\gamma}.
\]
当 \(\beta>1\) 时，\((\log u)^{|\gamma|}\) 的增长慢于任意正幂，故可与某个可积的 \(u^{-p}\)、\(p>1\) 比较。当 \(\beta<1\) 时，同样的比较给出发散。若 \(\beta=1\)，再令 \(v=\log u\)，便得到 \(\int v^{-\gamma}\,\mathrm dv\)，它收敛当且仅当 \(\gamma>1\)。

所求参数范围因而是
\[
 \boxed{\ \beta>1\quad\text{或}\quad\beta=1,\ \gamma>1\ }.
\]
必要性来自\cref{b3:thm:18-blaschke-zeros}。在这个范围内，按重数重复各 \(a_n\) 所得点列在圆盘内部无聚点，且 Blaschke 和有限。同一定理构造的乘积模不超过一、非恒零，并在每个 \(a_n\) 恰有 \(m_n\) 重零点，没有其他零点，从而证明充分性。
\end{solution}

% 自拟：从单因子误差导出乘积截断误差，并计算几何零点列的首阶尾乘积渐近。
\Needspace{6\baselineskip}
\begin{exercise}[Blaschke 尾乘积的误差]\label{b3:ex:18-draft-04}
设非零圆盘点列 \((a_j)_{j\ge1}\) 满足 Blaschke 条件，按重数重复，令
\[
 B_N=\prod_{j=1}^Nb_{a_j},\quad B=\prod_{j=1}^\infty b_{a_j},\quad
 T_N=\prod_{j>N}b_{a_j},\quad
 \delta_N=\sum_{j>N}(1-|a_j|).
\]
这里因子采用 \(b_a(0)=|a|\) 的归一化，\(B=B_NT_N\)。对 \(0<r<1\)，证明
\[
 \sup_{|z|\le r}|T_N(z)-1|
 \le\frac{1+r}{1-r}\delta_N,\quad
 \sup_{|z|\le r}|B(z)-B_N(z)|
 \le\frac{1+r}{1-r}\delta_N.
\]
进一步，若 \(a_j=1-q^j\)、\(0<q<1\)，证明在单位圆盘内局部一致地有
\[
 q^{-N}\bigl(T_N(z)-1\bigr)
 \longrightarrow-\frac{q}{1-q}\frac{1+z}{1-z}.
\]
\end{exercise}
\begin{solution}
对 \(M>N\)，有限乘积逐项相减给出
\[
 \prod_{j=N+1}^Mb_{a_j}-1
 =\sum_{k=N+1}^M(b_{a_k}-1)\prod_{j=N+1}^{k-1}b_{a_j}.
\]
每个因子在圆盘内模不超过一，再用正文的单因子误差界，得到
\[
 \sup_{|z|\le r}\left|\prod_{j=N+1}^Mb_{a_j}(z)-1\right|
 \le\frac{1+r}{1-r}\sum_{k=N+1}^M(1-|a_k|).
\]
令 \(M\to\infty\)，局部一致乘积收敛给出第一个估计。由 \(B-B_N=B_N(T_N-1)\) 及 \(|B_N|\le1\)，第二个估计随即成立。即使 \(B_N\) 在某点为零，乘积恒等式仍成立，无需在那里直接相除。

现在取 \(a_j=1-q^j\)。固定 \(r<1\)，记
\[
 h(z)=\frac{1+z}{1-z},\quad
 d_j(z)=1-b_{a_j}(z)
 =q^j\frac{1+z}{1-(1-q^j)z}.
\]
在 \(|z|\le r\) 上，分母与 \(1-z\) 均有共同正下界。将两分式相减可得
\[
 d_j(z)=q^jh(z)+O_r(q^{2j}),\quad |d_j(z)|=O_r(q^j),
\]
其中 \(q\) 固定，常数不依赖 \(j,z\)。当 \(N\) 足够大时，所有 \(j>N\) 都有 \(|d_j|\le1/2\)，故可以用原点附近的幂级数取 \(\log(1-d_j)\)，并一致地写成
\[
 \log(1-d_j)=-d_j+O_r(|d_j|^2).
\]
这些对数的级数绝对一致收敛，其指数正是 \(T_N\)。于是
\[
 \log T_N(z)
 =-h(z)\frac{q^{N+1}}{1-q}+O_r(q^{2N}).
\]
右侧在闭盘上一致为 \(O_r(q^N)\)，指数展开中的二次余项因而仍为 \(O_r(q^{2N})\)。所以
\[
 T_N(z)-1=-h(z)\frac{q^{N+1}}{1-q}+O_r(q^{2N}).
\]
乘以 \(q^{-N}\) 并令 \(N\to\infty\)，得到所求局部一致极限。
\end{solution}

% 自拟：在恒零极限中归一化零点质量，以圆周对数的绝对积分给出明确的 L1 误差。
\Needspace{6\baselineskip}
\begin{exercise}[趋零函数的归一化零点分布]\label{b3:ex:18-draft-05}
固定 \(0<r<1\)，令
\[
 f_n(z)=2^{-n}(z^n-r^n),\quad
 \nu_n=\frac1n\sum_{k=0}^{n-1}\delta_{r e^{2\pi\mathrm i k/n}},\quad
 u_n(z)=\frac1n\log|f_n(z)|.
\]
证明 \(f_n\) 在闭单位圆盘上一致趋于零，而 \(\nu_n\) 对紧支集连续测试函数收敛于半径 \(r\) 圆周上的均匀概率测度 \(\sigma_r\)。再证明
\[
 u_n\longrightarrow u,\quad
 u(z)=\log\max\{|z|,r\}-\log2
 \quad\text{于 }L^1_{\mathrm{loc}}(\mathbb D),
\]
并对每个 \(0<R<1\) 给出估计
\[
 \|u_n-u\|_{L^1(B(0,R))}\le\frac{2\pi R^2\log2}{n}.
\]
说明这些结论与正文的非零全纯极限条件有何关系。
\end{exercise}
\begin{solution}
在 \(|z|\le1\) 上，\(|f_n(z)|\le2^{-n}(1+r^n)\to0\)。方程 \(z^n=r^n\) 的根恰为题中列出的 \(n\) 个互异点，故 \(\nu_n\) 是零点计数测度除以 \(n\)。对 \(\varphi\in C_c(\mathbb D)\)，函数 \(t\mapsto\varphi(re^{\mathrm it})\) 连续且周期，其等分 Riemann 和给出
\[
 \int\varphi\,\mathrm d\nu_n
 =\frac1n\sum_{k=0}^{n-1}\varphi(re^{2\pi\mathrm i k/n})
 \longrightarrow\frac1{2\pi}\int_0^{2\pi}\varphi(re^{\mathrm it})\,\mathrm dt
 =\int\varphi\,\mathrm d\sigma_r.
\]

为控制对数奇性，先记 \(L_s(t)=\log|1-se^{\mathrm it}|\)，其中 \(0\le s\le1\)。圆周对数平均公式，包括 \(s=1\) 时的可积边界奇性，给出
\[
 \frac1{2\pi}\int_0^{2\pi}L_s(t)\,\mathrm dt=0.
\]
同时 \(L_s^+\le\log2\)。正、负部分积分相等，所以
\[
 \frac1{2\pi}\int_0^{2\pi}|L_s(t)|\,\mathrm dt
 =\frac1\pi\int_0^{2\pi}L_s^+(t)\,\mathrm dt
 \le2\log2.
\]
这是一条对全部 \(s\in[0,1]\) 有效的绝对积分界。

写 \(z=\rho e^{\mathrm i\theta}\)。当 \(0\le\rho\le r\) 时，从 \(z^n-r^n\) 中取出 \(-r^n\)；当 \(\rho>r\) 时取出 \(z^n\)。由此在几乎处处有
\[
 u_n(z)-u(z)=
 \begin{cases}
 n^{-1}\log|1-(\rho/r)^ne^{\mathrm in\theta}|,&0\le\rho\le r,\\
 n^{-1}\log|1-(r/\rho)^ne^{-\mathrm in\theta}|,&\rho>r.
 \end{cases}
\]
两种情形中的实参数都位于 \([0,1]\)。由于角函数为 \(2\pi\) 周期，把 \(\theta\) 换成 \(\pm n\theta\) 不改变一周平均；因此对每个固定 \(\rho\)，有
\[
 \int_0^{2\pi}|u_n(\rho e^{\mathrm i\theta})-u(\rho e^{\mathrm i\theta})|
 \,\mathrm d\theta\le\frac{4\pi\log2}{n}.
\]
对非负函数使用极坐标积分，得到
\[
 \int_{B(0,R)}|u_n-u|\,\mathrm dA
 \le\frac{4\pi\log2}{n}\int_0^R\rho\,\mathrm d\rho
 =\frac{2\pi R^2\log2}{n}.
\]
它也证明各个差值局部绝对可积，并给出所需 \(L^1\) 收敛。

分布公式还给 \((2\pi)^{-1}\Delta u_n=\nu_n\)，因而上述两个极限相容，并有 \((2\pi)^{-1}\Delta u=\sigma_r\)。这里全纯函数的极限是恒零函数，正文的对数模稳定性定理不适用；我们收敛的是另外除以 \(n\) 的对数模与零点质量。未归一化的零点总质量是 \(n\)，在一个等于一于该圆周附近的非负紧支集测试函数上趋于无穷，并没有趋向零测度。
\end{solution}

% 自拟：用高重数相邻零极点区分分布收敛与连续测试下的测度收敛，显式展示总变差失控。
\Needspace{6\baselineskip}
\begin{exercise}[零极点相消与测试函数]\label{b3:ex:18-draft-06}
对 \(n\ge2\)，令 \(\varepsilon_n=n^{-2}\)，并在单位圆盘考虑亚纯函数
\[
 F_n(z)=\left(\frac{z-\varepsilon_n}{z+\varepsilon_n}\right)^n,
 \quad
 \lambda_n=n\delta_{\varepsilon_n}-n\delta_{-\varepsilon_n}.
\]
证明 \(\log|F_n|\to0\) 于 \(L^1_{\mathrm{loc}}(\mathbb D)\)，且 \(\lambda_n\) 在分布意义下趋于零。构造一个 \(\varphi\in C_c(\mathbb D)\)，使 \(\int\varphi\,\mathrm d\lambda_n\) 不趋于零，并证明 \(\lambda_n\) 的正、负部分在原点的每个固定邻域内都没有共同有限质量上界。
\end{exercise}
\begin{solution}
对不在实线段 \([-\varepsilon,\varepsilon]\) 上的 \(z\)，沿实参数积分 \(\log|z-t|\) 的导数可得
\[
 \bigl|\log|z-\varepsilon|-\log|z+\varepsilon|\bigr|
 \le\int_{-\varepsilon}^{\varepsilon}\frac{\mathrm dt}{|z-t|}.
\]
排除的线段面积为零，不影响下面的面积积分。对 \(0<R<1\)、\(|t|\le1\)，平移后将积分区域放大为半径 \(R+1\) 的圆盘，得到
\[
 \int_{B(0,R)}\frac{\mathrm dA(z)}{|z-t|}
 \le\int_{B(0,R+1)}\frac{\mathrm dA(w)}{|w|}
 =2\pi(R+1).
\]
右侧核非负，可以用 Tonelli 定理换序。乘上 \(n\) 并取 \(\varepsilon=\varepsilon_n\)，便有
\[
 \int_{B(0,R)}|\log|F_n(z)||\,\mathrm dA(z)
 \le4\pi(R+1)n\varepsilon_n
 =\frac{4\pi(R+1)}n\longrightarrow0.
\]
每个内部紧集都包含于这样的圆盘，故局部 \(L^1\) 收敛成立。零点和极点处的无穷值按局部可积等价类理解。

由亚纯零极点公式，\((2\pi)^{-1}\Delta\log|F_n|=\lambda_n\)，局部 \(L^1\) 收敛已经推出分布收敛。也可以直接检查：对 \(\psi\in C_c^\infty(\mathbb D)\)，当 \(n\) 足够大时，
\[
 \left|\int\psi\,\mathrm d\lambda_n\right|
 =n|\psi(\varepsilon_n)-\psi(-\varepsilon_n)|
 \le2n\varepsilon_n\sup|\partial_x\psi|
 \longrightarrow0.
\]
这个估计用到了光滑测试函数的一阶变化控制。

现在取非负 \(\chi\in C_c^\infty(\mathbb D)\)，使其在原点附近恒等于一，定义
\[
 \varphi(x+\mathrm iy)=\chi(x+\mathrm iy)\operatorname{sgn}(x)\sqrt{|x|},
 \quad \operatorname{sgn}(0)\sqrt0=0.
\]
该函数连续且有紧支集，但在 \(x=0\) 附近不满足前面的一阶变化估计。对充分大的 \(n\)，有
\[
 \int\varphi\,\mathrm d\lambda_n
 =n\left(\sqrt{\varepsilon_n}-(-\sqrt{\varepsilon_n})\right)=2.
\]
所以光滑测试下的分布收敛并没有推广为全部紧支集连续测试下的测度收敛。

两个支点互异，故 Jordan 分解为
\[
 \lambda_n^+=n\delta_{\varepsilon_n},\quad
 \lambda_n^-=n\delta_{-\varepsilon_n}.
\]
在任何固定原点邻域内，两点最终都被包含，正、负质量各为 \(n\)。特别地，对前面的 \(\chi\)，最终有
\(\int\chi\,\mathrm d\lambda_n^+=\int\chi\,\mathrm d\lambda_n^-=n\)，总变差积分为 \(2n\)。缺少局部总变差上界，便不能把连续函数的光滑逼近误差统一控制；这是两类测试得出不同结论的具体原因。
\end{solution}

% 对原自拟题实质改编：在有限零点除法后加入两个插值条件，以 Schwarz 引理给出可解判据和刚性。
\Needspace{6\baselineskip}
\begin{exercise}[指定零点后的两点插值]\label{b3:ex:18-draft-07}
设 \(B\) 为单位圆盘内的有限 Blaschke 乘积，\(p,q\in\mathbb D\) 互异且 \(B(p)B(q)\ne0\)，给定 \(c,d\in\mathbb C\)。求存在圆盘全纯函数 \(f\) 满足
\[
 |f|\le1,\quad f/B\text{ 在圆盘全纯},\quad f(p)=c,\quad f(q)=d
\]
的充要条件。商条件表示 \(f\) 在指定各点至少具有 \(B\) 的相应零点重数；允许额外零点，也允许 \(f\equiv0\)。记
\[
 A=\frac{c}{B(p)},\quad C=\frac{d}{B(q)},\quad
 T_a(z)=\frac{z-a}{1-\overline a z}\quad(a\in\mathbb D).
\]
分别讨论 \(|A|>1\)、\(|A|=1\) 和 \(|A|<1\)。在最后一种情形下，若所得的 Schwarz 距离不等式取等号，确定解的形式及唯一性。
\end{exercise}
\begin{solution}
设已有这样的 \(f\)，令 \(h=f/B\)。有限 Blaschke 乘积在单位圆周上模为一，并连续到圆周。在包住其全部零点的圆周 \(|z|=s<1\) 上，最大模原理给
\[
 |h(z)|\le\frac1{\min_{|\zeta|=s}|B(\zeta)|}\quad(|z|<s).
\]
令 \(s\uparrow1\)，可得 \(|h|\le1\)。反过来，任意这样的 \(h\) 都使 \(f=Bh\) 满足模界和零点要求。因此问题等价于求圆盘全纯函数 \(h\)，使 \(|h|\le1\)、\(h(p)=A\)、\(h(q)=C\)。

若 \(|A|>1\)，无解。若 \(|A|=1\)，最大模原理强制 \(h\equiv A\)，所以有解当且仅当 \(C=A\)，此时唯一解为 \(f=AB\)。

若 \(|A|<1\)，任何可行 \(h\) 都将圆盘映入开圆盘；否则它在内部取到模一值，便应为模一常数，与 \(h(p)=A\) 矛盾。因而必须有 \(C\in\mathbb D\)。函数
\[
 g=T_A\circ h\circ T_p^{-1}
\]
将圆盘映入自身且 \(g(0)=0\)。Schwarz 引理在点 \(T_p(q)\ne0\) 给出必要条件
\[
 |T_A(C)|\le|T_p(q)|.
\]
反之，若 \(C\in\mathbb D\) 且该不等式成立，令
\[
 \lambda=\frac{T_A(C)}{T_p(q)},\quad
 h(z)=T_A^{-1}\bigl(\lambda T_p(z)\bigr).
\]
因为 \(|\lambda|\le1\)，这个 \(h\) 在圆盘全纯，取值在圆盘内，并满足两个插值条件。所以完整判据为
\[
 \begin{cases}
 \text{无解},&|A|>1,\\
 C=A,&|A|=1,\\
 C\in\mathbb D\ \text{且 }|T_A(C)|\le|T_p(q)|,&|A|<1.
 \end{cases}
\]

在 \(|A|<1\) 的情形，若不等式取等号，Schwarz 引理的等号结论使 \(g(\zeta)=\lambda\zeta\)、\(|\lambda|=1\)。两个插值值已唯一决定 \(\lambda\)，所以唯一的解就是
\[
 f(z)=B(z)T_A^{-1}\bigl(\lambda T_p(z)\bigr).
\]
后一因子为圆盘自同构，恰增加一个简单零点；它也可能落在 \(B\) 的既有零点处，使该点的总重数增加一。题设只要求保留指定零点及至少那些重数，因此这两种情况都允许。
\end{solution}

% 对原自拟题实质改编：从倒反不变性推进到几何平均的双侧系数界、最优常数和等号判别。
\Needspace{6\baselineskip}
\begin{exercise}[多项式的几何平均与系数界]\label{b3:ex:18-draft-08}
设非零多项式 \(P(z)=\sum_{k=0}^dc_kz^k\) 的次数为 \(d\)，定义
\[
 \mathcal M(P)=\exp\left(\frac1{2\pi}\int_0^{2\pi}
 \log|P(e^{\mathrm it})|\,\mathrm dt\right),\quad
 P^*(z)=z^d\overline{P(1/\overline z)},
\]
其中 \(P^*\) 在原点按多项式延拓。将 \(P\) 的根按重数列为 \(a_1,\ldots,a_d\)，证明
\[
 \mathcal M(P)=|c_d|\prod_{j=1}^d\max\{1,|a_j|\},\quad
 \mathcal M(P^*)=\mathcal M(P),
\]
以及
\[
 \max_{0\le k\le d}\frac{|c_k|}{\binom dk}
 \le\mathcal M(P)
 \le\left(\sum_{k=0}^d|c_k|^2\right)^{1/2}.
\]
证明每个系数不等式中的常数 \(\binom dk\) 都不能减小，并确定右端不等式取等号的全部多项式。允许原点零点与单位圆周上的零点。
\end{exercise}
\begin{solution}
因式分解给 \(P(z)=c_d\prod_{j=1}^d(z-a_j)\)。每个因子的圆周对数都绝对可积，平均为 \(\log\max\{1,|a_j|\}\)，其中单位圆周上的对数奇性也包括在内。逐项相加并指数化，得到根的乘积公式。\(d=0\) 时空乘积取一。

展开倒反定义可得
\[
 P^*(z)=\sum_{k=0}^d\overline{c_k}z^{d-k}
 =\overline{c_d}\prod_{j=1}^d(1-\overline{a_j}z).
\]
这是非零多项式；若 \(c_0=0\)，它的实际次数可以小于 \(d\)，但定义仍有效。在单位圆周上 \(|P^*|=|P|\)，两者的可积对数具有同一平均，所以 \(\mathcal M(P^*)=\mathcal M(P)\)。

系数 \(c_k\) 是 \(c_d\) 乘以 \(d-k\) 阶根的基本对称和，再乘一个符号。每个由 \(d-k\) 个根组成的乘积都满足
\[
 \prod_{j\in I}|a_j|\le\prod_{j=1}^d\max\{1,|a_j|\}.
\]
这样的项有 \(\binom d{d-k}=\binom dk\) 个，三角不等式便给
\( |c_k|\le\binom dk\mathcal M(P)\)。取全部 \(k\) 的最大值得到左界。若 \(P(z)=C(z-a)^d\)、\(C\ne0\)、\(|a|=1\)，则 \(\mathcal M(P)=|C|\)，而每个 \(|c_k|=|C|\binom dk\)。因此每个系数的常数都已最优。

为证右界，记
\[
 X(t)=|P(e^{\mathrm it})|^2,\quad
 A=\frac1{2\pi}\int_0^{2\pi}X(t)\,\mathrm dt>0.
\]
由 \(\log x\le x-1\)，在 \(X>0\) 的点有
\(\log(X/A)\le X/A-1\)。有限个圆周零点不影响积分，而 \(\log X\) 已知可积，故
\[
 2\log\mathcal M(P)-\log A
 \le\frac1{2\pi}\int_0^{2\pi}(X/A-1)\,\mathrm dt=0.
\]
再将有限 Fourier 和展开积分，异频项的积分为零，得到
\[
 A=\frac1{2\pi}\int_0^{2\pi}
 \sum_{j,k=0}^dc_j\overline{c_k}e^{\mathrm i(j-k)t}\,\mathrm dt
 =\sum_{k=0}^d|c_k|^2.
\]
这证明了右界。

等号要求 \(X/A-1-\log(X/A)\) 几乎处处为零，即 \(X=A\) 几乎处处。由连续性，\(|P|\) 在整个单位圆周为常数。设 \(m\) 是最小的非零系数指标。若 \(m<d\)，则 \(|P(e^{\mathrm it})|^2\) 的频率 \(d-m\) 项只有一项，其系数为 \(c_d\overline{c_m}\ne0\)，与常数函数矛盾。所以 \(m=d\)，\(P(z)=c_dz^d\)。反过来，任一非零单项式都在单位圆周上有恒定模，确实取等号；其中包括非零常数。
\end{solution}
